\section{Model, symmetries, and topological classification}
\label{sec:model}

\subsection{Dimerized XXZ chain}

We study an open chain of $L$ spin-$1/2$ sites, with $X_j,Y_j,Z_j$ denoting Pauli operators and lattice spacing set to one:
\begin{equation}
H(s,\delta)=\sum_{j=1}^{L-1}t_j
\left[e^{-\delta}(X_jX_{j+1}+Y_jY_{j+1})
+e^\delta Z_jZ_{j+1}\right],
\label{eq:H}
\end{equation}
where $t_{2r-1}=1-s$, $t_{2r}=s$, and $0\leq s\leq1$.
The alternating bonds inherit the pattern of the SSH model~\cite{su1979solitons}.
The free-fermion band classification~\cite{ryu2010topological} provides the noninteracting reference.
Here the spin degrees of freedom and longitudinal interaction lead to a many-body symmetry classification.

To compare conventions with the extended bosonic SSH model of Ref.~\cite{elben2020manybody}, its odd/even couplings are $J=2(1-s)e^{-\delta}$ and $J'=2se^{-\delta}$, and its anisotropy is
\begin{equation}
\Delta=\frac{J_z}{J_{xx}}=e^{2\delta}.
\label{eq:anisotropy}
\end{equation}
Thus $\delta=0$ is the isotropic point, $\delta<0$ is easy plane, and $\delta>0$ is easy axis. This correspondence also connects our model to the interacting bosonic SSH chain realized with Rydberg atoms~\cite{deleseleuc2019observation}. Unless stated otherwise, $L=8$.

\subsection{Symmetries and the variational sector}

The Hamiltonian commutes with the total magnetization $Z_{\rm tot}=\sum_jZ_j$, global spin flip $G_x=\prod_jX_j$, and bond-centered reflection $R:j\mapsto L+1-j$.
It also respects spin time reversal. Continuous spin rotation is $U(1)$ for $\delta\neq0$ and enlarges to $SU(2)$ at $\delta=0$.
For the sizes used here, $L$ is a multiple of four and all three reference states have $Z_{\rm tot}=0$, $G_x=+1$, and $R=+1$.
The conserved operator
\begin{equation}
P=Z_{\rm tot}^2-G_x-R
\end{equation}
packages these conditions: $P=-2$ is equivalent to their simultaneous satisfaction. The ideal variational circuit preserves this sector.

\subsection{Why the two dimerizations are distinct}

A symmetry-protected topological (SPT) phase is defined by equivalence under symmetry-preserving local transformations of gapped states~\cite{gu2009tensor,chen2010local}.
For the onsite symmetries of this chain, the natural unit cell is $(2r-1,2r)$.
Its total spin is $T_r^a=(A_{2r-1}+A_{2r})/2$, with $A=X,Y,Z$ for $a=x,y,z$.
The cell Hilbert space contains both spin-zero and spin-one sectors.
The cell rotations $u_x=e^{i\pi T_r^x}$ and $u_z=e^{i\pi T_r^z}$ commute and generate $D_2\simeq\mathbb Z_2\times\mathbb Z_2$.

For a symmetric, injective matrix product state (MPS), the symmetry acts on its virtual bonds as
\begin{equation}
\sum_n(u_g)_{mn}A^n=e^{i\phi_g}V_g^\dagger A^mV_g,
\qquad V_gV_h=\omega(g,h)V_{gh}.
\label{eq:mps-symmetry}
\end{equation}
The inequivalent projective representations, labeled by $H^2(G,U(1))$ for unitary onsite $G$, distinguish the symmetric gapped phases~\cite{chen2011classification,schuch2011classifying}.
The MPS and Schmidt-decomposition framework is reviewed in Refs.~\cite{schollwock2011density,cirac2021matrix}.
For $D_2$, the gauge-invariant relation
\begin{equation}
V_xV_z=\eta V_zV_x,\qquad \eta=\pm1
\label{eq:projective-index}
\end{equation}
separates the trivial and Haldane-type classes~\cite{pollmann2012detection}.

At $s=0$, singlets lie within cells, giving a product of cell states and $\eta=+1$.
At $s=1$, bulk singlets cross cell boundaries, leaving spin-$1/2$ edge degrees of freedom with $\eta=-1$.
Our topological reference closes the two edge spins into a singlet; this selects one state of the open-chain edge manifold.
This dimer limit illustrates symmetry fractionalization~\cite{verresen2017onedimensional}.
Bond inversion or time reversal can independently protect the Haldane-type distinction~\cite{pollmann2012symmetry}.
The entanglement spectrum, introduced in Ref.~\cite{li2008entanglement}, has symmetry-protected degeneracies in this one-dimensional setting~\cite{pollmann2010entanglement}.
These classification statements concern gapped bulk phases without spontaneous breaking of the protecting symmetry. The AFM region instead has broken spin-flip symmetry in the thermodynamic limit; a finite-chain symmetric cat state can retain the symmetry while displaying large staggered correlations.

\subsection{Partial-reflection diagnostic}

We adopt the bosonic partial-reflection diagnostic of Ref.~\cite{elben2020manybody}.
Partial spatial transformations also enter constructions of many-body invariants for fermionic SPT phases~\cite{shiozaki2017manybody}.
For a central interval $I=I_1\cup I_2$, let $\rho_I$ be its reduced density matrix and $\mathcal R_I$ reverse its sites. We evaluate
\begin{equation}
\tilde Z_{\mathcal R}
=\frac{\operatorname{Tr}(\rho_I\mathcal R_I)}
{\sqrt{\left[\operatorname{Tr}(\rho_{I_1}^2)
+\operatorname{Tr}(\rho_{I_2}^2)\right]/2}}.
\label{eq:reflection-marker}
\end{equation}
For $L=8$, $I=(3,4,5,6)$ and its halves are $(3,4)$ and $(5,6)$.
In the symmetric gapped regimes, the quantized limits $+1$ and $-1$ require intervals large compared with the correlation length; the AFM limit is zero in the corresponding model~\cite{elben2020manybody}.
We use the four-site interval to evaluate the finite-size phase map.

We compute Eq.~\eqref{eq:reflection-marker} from numerical wave functions.
The expression combines off-diagonal matrix elements of $\mathcal R_I$ with subsystem purities.
Randomized measurements with suitable basis rotations provide an experimental route to these quantities~\cite{elben2022toolbox}.
